\section{Example II: electrodynamics, special relativity and the point particle}
\label{sec:sr}

\subsection{Lorentz covariance of Maxwell's equations}

With $F^{0i}=-E^i$ and $F^{ij}=-\epsilon^{ijk}B^k$, the Maxwell equations in
Gaussian units read
\begin{equation}\label{eq:maxwell}
  \partial_\mu F^{\mu\nu}=\frac{4\pi}{c}\,j^\nu,\qquad
  \partial_\mu\tilde F^{\mu\nu}=0,\qquad
  \tilde F^{\mu\nu}=\tfrac12\epsilon^{\mu\nu\alpha\beta}F_{\alpha\beta}.
\end{equation}
Under a Lorentz transformation \eqref{eq:poincare}, $F^{\mu\nu}$ transforms as a
rank-two tensor and $j^\nu$ as a vector, so
$\mathcal E^\nu\equiv\partial_\mu F^{\mu\nu}-\tfrac{4\pi}{c}j^\nu$ is a
four-vector: $\mathcal E'^\nu=\Lambda^\nu{}_\alpha\mathcal E^\alpha$. This is
\eqref{eq:covariant-eq} with $D_E(g)=\Lambda$; the metric $\eta_{\mu\nu}$,
needed to raise and lower indices, is invariant, so Proposition~\ref{prop:main}
gives Poincar\'e symmetry.

It is worth stating exactly what is covariant and what is invariant.
The fields ${\bm{E}}$ and ${\bm{B}}$ are \emph{components} of a tensor and mix under
boosts; for a boost with velocity $c{\bm{\beta}}$,
\begin{equation}\label{eq:boost-fields}
  {\bm{E}}'_\parallel={\bm{E}}_\parallel,\quad {\bm{B}}'_\parallel={\bm{B}}_\parallel,\quad
  {\bm{E}}'_\perp=\gamma({\bm{E}}_\perp+{\bm{\beta}}\times{\bm{B}}),\quad
  {\bm{B}}'_\perp=\gamma({\bm{B}}_\perp-{\bm{\beta}}\times{\bm{E}}).
\end{equation}
The \emph{invariants} are scalars built from $F$ and $\eta$ (and $\epsilon$):
\begin{equation}\label{eq:EM-invariants}
  F_{\mu\nu}F^{\mu\nu}=2(B^2-E^2),\qquad
  F_{\mu\nu}\tilde F^{\mu\nu}=-4\,{\bm{E}}\!\cdot\!{\bm{B}} .
\end{equation}
A purely electric field acquires a magnetic component in a moving frame, yet
$E^2-B^2$ and ${\bm{E}}\!\cdot\!{\bm{B}}$ do not change. For example, a field with
${\bm{E}}\!\cdot\!{\bm{B}}=0$ and $E^2>B^2$ can be boosted to a frame in which it is
purely electric, but never to one in which it is purely magnetic. The
Lagrangian density in \eqref{eq:maxwell-action} is $\tfrac1{8\pi}(E^2-B^2)$, a
Lorentz scalar: covariance of the action in the sense of
Section~\ref{sec:definitions}. (Because $\epsilon^{\mu\nu\alpha\beta}$ is a
pseudotensor, covariance under parity requires $\tilde F$ to be treated as
one.)

The same equations are \emph{not} covariant under Galilean transformations
with any local rule for ${\bm{E}},{\bm{B}}$; limits in which they are (``electric'' and
``magnetic'' Galilean electromagnetism) exist, but each captures only part of
the equations \cite{lbll}. The failure of Galilean covariance of Maxwell's
theory was the original motivation for special relativity.

\subsection{Larger symmetries and symmetry without an action}

The Poincar\'e group is not the whole story.
\begin{itemize}[leftmargin=1.6em,itemsep=2pt]
\item \emph{Conformal symmetry.} In four dimensions the sourceless Maxwell
  equations are covariant under all diffeomorphisms that preserve the
  \emph{conformal class} $[\eta]$, i.e.\ $\eta\to\Omega^2\eta$; the
  corresponding symmetry group is the 15-parameter conformal group
  (translations 4, Lorentz 6, dilatation 1, special conformal 4). The
  stabilizer of $[\eta]$ is bigger than the stabilizer of $\eta$, which is why
  the symmetry is bigger: Proposition~\ref{prop:main} with a weaker background.
\item \emph{Duality.} In the absence of sources, the replacement
  ${\bm{E}}\to{\bm{E}}\cos\alpha+{\bm{B}}\sin\alpha$, ${\bm{B}}\to-{\bm{E}}\sin\alpha+{\bm{B}}\cos\alpha$
  maps solutions to solutions, so it is a symmetry of the \emph{equations}.
  It is not a local variational symmetry of the action \eqref{eq:maxwell-action}
  built from $A_\mu$ (for $\alpha=\pi/2$ the Lagrangian density changes sign),
  though conserved charges can be obtained in a reformulation involving a
  second potential \cite{deser}. This illustrates the failure of the converse
  in Proposition~\ref{prop:variational}.
\end{itemize}

\subsection{The relativistic point particle: reparametrization as gauge symmetry}

For a free particle, the manifestly covariant action is
\begin{equation}\label{eq:particle}
  S=-mc\int\sqrt{\eta_{\mu\nu}\,\dot x^\mu\dot x^\nu}\;\dd\lambda ,
  \qquad \dot{}\equiv\frac{\dd}{\dd\lambda}.
\end{equation}
It is invariant under two very different kinds of transformation.

\emph{(i) Poincar\'e transformations} \eqref{eq:poincare}, a global
symmetry (the Lagrangian is a scalar and $\eta$ is invariant). Noether's first
theorem gives the conserved quantities $\pi_\mu=\partial L/\partial\dot x^\mu$
and $J^{\mu\nu}=x^\mu\pi^\nu-x^\nu\pi^\mu$.

\emph{(ii) Reparametrizations} $\lambda\to f(\lambda)$, which depend on an
arbitrary function: a local symmetry, $\delta x^\mu=\eps(\lambda)\dot x^\mu$.
Noether's second theorem \eqref{eq:noether2} gives the identity
\begin{equation}\label{eq:param-identity}
  \dot x^\mu\,\mathrm{EL}_\mu\equiv0,
\end{equation}
so only three of the four equations for $x^\mu(\lambda)$ are independent, and
one function (the parametrization) is undetermined. In the Hamiltonian
description the momenta satisfy the identity
$\eta^{\mu\nu}\pi_\mu\pi_\nu-m^2c^2=0$ (a primary constraint), the canonical
Hamiltonian vanishes because $L$ is homogeneous of degree one in the
velocities, and the dynamics is generated by the constraint times an arbitrary
multiplier.

Choosing the gauge $\lambda=t$ gives
$L=-mc^2\sqrt{1-v^2/c^2}$, which is \emph{not} manifestly covariant: the
boost acts \emph{nonlinearly} on $x^i(t)$, since a boost mixes $t$ with the
coordinates and the gauge condition must be restored afterwards. The
Poincar\'e symmetry is still there; it is merely hidden by the gauge choice.
Covariance is a property of the formulation, symmetry is not.

% ---------------------------------------------------------------------
\section{Example III: general relativity}
\label{sec:gr}

\subsection{General covariance is not a symmetry of the world}

Einstein's field equations in the convention of Landau and Lifshitz
(signature $(+,-,-,-)$),
\begin{equation}\label{eq:einstein}
  R_{\mu\nu}-\tfrac12 g_{\mu\nu}R=\frac{8\pi G}{c^4}\,T_{\mu\nu},
\end{equation}
derive from $S=-\tfrac{c^3}{16\pi G}\int R\sqrt{-g}\,\dd^4x+S_{\rm matter}$
and are covariant under the full diffeomorphism group $\diff(M)$. In
Proposition~\ref{prop:main}, $\mathcal G=\diff(M)$ and the metric $g_{\mu\nu}$
is \emph{dynamical}: there is no fixed background $b_0$ (in the absence of
fixed matter fields). The stabilizer is therefore the entire group,
\begin{equation}
  G=\stab_{\diff(M)}(\varnothing)=\diff(M),
\end{equation}
and by Remark~\ref{rem:gauge} diffeomorphisms are a \emph{gauge redundancy}
rather than a physical symmetry. Here lies the resolution of Kretschmann's
objection: Einstein's theory is not special because its equations can be
written covariantly (every theory can be), but because they contain
\emph{no absolute objects} \cite{anderson} (Section~\ref{sec:kretschmann} gives
the details). In special relativity written in arbitrary coordinates the metric $\eta$ is an absolute object; in GR it is
replaced by the dynamical $g_{\mu\nu}$.

\subsection{The hole argument}

Let $g$ be a solution on $M$, and let $\psi$ be a diffeomorphism equal to the
identity outside a compact region $H$. Then $\psi^*g$ is also a solution,
equal to $g$ outside $H$ and different inside. If points of $M$ were taken to
have physical identity (a metric field being \emph{the} object of physical
interest), determinism would fail: the Einstein equations plus data outside $H$
do not determine the field inside. The way out is to regard $g$ and $\psi^*g$
as one and the same physical situation. Observables are then
diffeomorphism-\emph{invariant} (relational) quantities: for instance,
proper times along worldlines between their intersections. This is exactly
the gauge-redundancy interpretation of Remark~\ref{rem:gauge}
\cite{earman,norton}.

\subsection{Noether's second theorem in GR: conservation only for Killing vectors}

Under $\delta g_{\mu\nu}=-(\nabla_\mu\xi_\nu+\nabla_\nu\xi_\mu)$ with arbitrary
$\xi^\mu(x)$ the action is invariant, and \eqref{eq:noether2} becomes the
contracted Bianchi identity $\nabla_\mu G^{\mu\nu}\equiv0$. Consequently the
field equations \eqref{eq:einstein} imply $\nabla_\mu T^{\mu\nu}=0$, which
\emph{is not a conservation law} in the sense of a conserved charge: in a
curved spacetime $\nabla_\mu T^{\mu\nu}=0$ does not integrate to a conserved
quantity. A conserved current is obtained only if there exists a Killing vector
$\xi$ of the particular solution under consideration:
\begin{equation}
  J^\mu=T^{\mu\nu}\xi_\nu,\qquad
  \nabla_\mu J^\mu=(\nabla_\mu T^{\mu\nu})\xi_\nu+T^{\mu\nu}\nabla_{(\mu}\xi_{\nu)}=0.
\end{equation}
This is the distinction of Section~\ref{sec:law-solution} in its
sharpest form: the \emph{law} has the symmetry $\diff(M)$ (gauge), but
conservation of energy and momentum rests on the symmetry of the \emph{solution}
(the isotropy group $G_{\phi_0}$ of the background spacetime, such as time
translation invariance of the Schwarzschild solution).

% ---------------------------------------------------------------------
\section{Kretschmann's observation in detail}
\label{sec:kretschmann}

Remark~\ref{rem:kretschmann} stated Kretschmann's observation in one
sentence. The confusion it addresses is among the most persistent in the
foundations of relativity, so we now examine it more closely: the historical
dispute, a constructive proof, what the proof does and does not show, and what
remains of ``general covariance'' afterwards.

\subsection{The dispute}

In 1915--16 Einstein presented general covariance, the form-invariance of the
field equations under arbitrary coordinate changes, as a physical principle,
for two reasons. First, it seemed to extend the principle of relativity from
inertial frames to arbitrary states of motion. Second, he argued that the
physical content of any theory is exhausted by \emph{coincidences}: the
intersection of two worldlines, the meeting of a pointer and a mark on a
scale. Coordinates merely label such coincidences, no labeling can be
physically privileged, and a theory that singles out a class of coordinates
would draw a distinction without empirical basis (the \emph{point-coincidence
argument}) \cite{norton}.

Kretschmann accepted the premise and reversed the conclusion
\cite{kretschmann}. If only coincidences are physical, he argued, then
\emph{any} spacetime theory, however un-relativistic, can be restated in a form
that does not depend on the labeling, at the price of mathematical effort: one
adds to the formalism whatever is needed to keep track of the labeling. The
demand for general covariance therefore tests the ingenuity of the theorist and
places no restriction on nature. Einstein conceded the logical point in 1918
\cite{einstein1918} but maintained that the demand retains a \emph{heuristic}
value: among generally covariant equations only the simplest and most natural
are plausible as fundamental laws, and the covariant versions of older theories
are comparatively baroque. Historically the dispute coincided with the work of
Hilbert, Klein and Noether on identities among the field equations of
generally covariant theories, and the second theorem of
Section~\ref{sec:noether} came out of that work \cite{ks}. Anderson and
Friedman later turned Kretschmann's observation into a precise criterion
\cite{anderson,friedman}, which we met in Remark~\ref{rem:smaller}; Norton's
survey \cite{norton} and his later discussion \cite{bc} are the best guides to
the debate.

\subsection{A constructive proof}

Let $E[\varphi]=0$ be \emph{any} system of differential equations for fields
$\varphi$, written in a fixed ``preferred'' chart $x^a$ on a region of
spacetime. It may be Galilean, Lorentzian, or covariant under no group at all.
In an arbitrary chart $y^\mu$ the preferred coordinates appear as four scalar
functions $X^a(y)$. Put
\begin{equation}\label{eq:soldering}
  e^a{}_\mu=\partial_\mu X^a,\qquad
  e_a{}^\mu\,e^b{}_\mu=\delta_a^b .
\end{equation}
The \emph{frame components} of a tensor field are its contractions with
$e_a{}^\mu$ and $e^a{}_\mu$ (for a covector, $\varphi_a=e_a{}^\mu\varphi_\mu$).
They carry only the index $a$ and are therefore \emph{scalars} under changes
of the chart $y$. The operators $D_a=e_a{}^\mu\partial_\mu$ map scalars to
scalars. Let $\tilde E[\tilde\varphi;e]$ be obtained from $E$ by replacing every
field by its frame components and every $\partial/\partial x^a$ by $D_a$.

\begin{proposition}[Kretschmann]\label{prop:kretschmann}
The system $\tilde E[\tilde\varphi;e]=0$ is $\diff(M)$-covariant in the sense
of \eqref{eq:covariant-eq}, with the soldering form $e^a{}_\mu$ (Section~\ref{sec:soldering}) as background.
It is equivalent to the original system: $(\tilde\varphi,e)$ solves
$\tilde E=0$ if and only if the frame components of $\tilde\varphi$ solve
$E=0$.
\end{proposition}
\begin{proof}
Each frame component and each $D_a(\cdot)$ is a scalar under $\diff(M)$ acting
on $(\tilde\varphi,e)$ by push-forward, so $\tilde E$ is built from scalars
and $\tilde E[\psi\!\cdot\!\tilde\varphi;\psi\!\cdot\! e]
=\psi\!\cdot\!\tilde E[\tilde\varphi;e]$ for every diffeomorphism $\psi$. In the
particular chart $y=X$ one has $e^a{}_\mu=\delta^a_\mu$ and $D_a=\partial/\partial x^a$,
so $\tilde E=E$. \qedhere
\end{proof}

Nothing about $E$ was used, which is precisely the content of Kretschmann's
claim: general covariance is available for any theory, at the cost of a
background field $e^a{}_\mu$ that the original formulation did not display.

\begin{example}[Wave equation]\label{ex:wave}
For $\eta^{ab}\partial_a\partial_b\varphi=0$ the construction gives
$\eta^{ab}D_aD_b\varphi=0$. Expanding $D_aD_b\varphi=e_a{}^\mu\partial_\mu(e_b{}^\nu\partial_\nu\varphi)$
and using $\eta^{ab}e_a{}^\mu e_b{}^\nu=\eta^{\mu\nu}$,
\begin{equation}
  \eta^{ab}D_aD_b\varphi
  =\eta^{\mu\nu}\big(\partial_\mu\partial_\nu\varphi-\Gamma^\lambda{}_{\mu\nu}\partial_\lambda\varphi\big)
  =\eta^{\mu\nu}\nabla_\mu\partial_\nu\varphi,\qquad
  \eta_{\mu\nu}=\eta_{ab}e^a{}_\mu e^b{}_\nu,
\end{equation}
where $\Gamma^\lambda{}_{\mu\nu}=e_a{}^\lambda\partial_\mu e^a{}_\nu$ are the
Christoffel symbols of $\eta_{\mu\nu}$ (flat, since $e^a{}_\mu$ is a gradient). This is exactly
how Christoffel symbols enter the ``curvilinear'' form of every equation of
special relativity.
\end{example}

\begin{example}[Newtonian mechanics, parametrized]\label{ex:param}
The free Newtonian particle, $S=\int\tfrac12m\,\dd{\bm{x}}^2/\dd t$, becomes
generally covariant (here: reparametrization invariant) if $t$ is promoted to a
dynamical variable of the worldline,
\begin{equation}
  S=\int\frac m2\,\frac{\dot{{\bm{x}}}^2}{\dot t}\,\dd\lambda,\qquad
  \dot{}=\frac{\dd}{\dd\lambda}.
\end{equation}
This is homogeneous of degree one in $(\dot t,\dot{\bm{x}})$, hence invariant under
$\lambda\to f(\lambda)$; choosing $\lambda=t$ gives back the original
action. The momenta obey $p_t+{\bm{p}}^2/2m\approx0$, the Newtonian analogue of the
mass-shell constraint of Section~\ref{sec:sr}, with $-p_t$ playing the
role of the energy. This is Kretschmann's observation in one dimension.
\end{example}

\begin{remark}[Dependence on the soldering form]\label{rem:soldering-dep}
If $E$ is covariant under a group $G_0$ acting on the frame index $a$ (Lorentz
transformations, say), then $\tilde E$ depends on $e^a{}_\mu$ only through its
$G_0$-invariant parts: for the Poincar\'e group through the flat metric
$\eta_{\mu\nu}=\eta_{ab}e^a{}_\mu e^b{}_\nu$ alone, and for the Galilei group
through the Newton--Cartan data of Section~\ref{sec:galilei}.
\end{remark}

\subsection{The soldering form}
\label{sec:soldering}

The original formulation of $E$ contains two kinds of index that were never
distinguished, because they are numerically identified in the preferred chart.
The index $a$ labels the coordinates of the \emph{model space} $V=\mathbb R^4$,
on which the group $G_0$ (Poincar\'e, Galilei, \dots) acts and in which the
original theory lives. The index $\mu$ labels directions at a point of the
spacetime manifold $M$, whose charts $y^\mu$ are arbitrary. Covariantization
must pull the two apart: $\mu$ is acted on by $\diff(M)$, whereas $a$ is merely
a label, acted on at most by $G_0$.

The one-form $e^a{}_\mu\,\dd y^\mu=\dd X^a$ is the link between them. It is a
one-form on $M$ with values in $V$: at each point $p$ it defines the linear map
\begin{equation}
  e_p\colon T_pM\to V,\qquad v^\mu\mapsto e^a{}_\mu\,v^\mu ,
\end{equation}
which is invertible precisely because $X^a$ is a chart. It identifies every
tangent space with the model space, point by point, and thereby \emph{solders}
(welds) the model space onto the manifold; hence the name. In the mathematical
literature the \emph{solder form} (or canonical form) is the tautological
$V$-valued one-form on the frame bundle, and its pull-backs by local sections
are the coframes, tetrads or vierbeins $e^a{}_\mu$ \cite{kn}. In Cartan's
generalization of Klein's Erlangen programme, a geometry is a manifold whose
tangent spaces are soldered to the model space of a Klein geometry $G_0/H$
\cite{sharpe}. (Some authors reserve ``soldering form'' for the canonical
form and call $e^a{}_\mu$ a tetrad; we follow the usage of physics.)

The soldering form does three jobs in the construction above.
\begin{enumerate}[leftmargin=1.8em,itemsep=3pt,label=\textbf{\arabic*.}]
\item \emph{It converts indices.} Frame components $\varphi_a=e_a{}^\mu\varphi_\mu$
  and the derivatives $D_a=e_a{}^\mu\partial_\mu$ are scalars under
  $\diff(M)$ because $e^a{}_\mu$ is the only place where an index of each kind
  meets.
\item \emph{It records the preferred chart and nothing else.} Since
  $e^a{}_\mu=\partial_\mu X^a$, it is closed, $\dd e^a=0$. A translation
  $X^a\to X^a+c^a$ leaves $e^a{}_\mu$ unchanged, and a Lorentz transformation
  $X^a\to\Lambda^a{}_bX^b$ rotates it, $e\to\Lambda e$. Only the
  $G_0$-invariant content matters (Remark~\ref{rem:soldering-dep}): for the
  Poincar\'e group the metric $\eta_{\mu\nu}=\eta_{ab}e^a{}_\mu e^b{}_\nu$; for
  the Galilei group the time one-form $t_\mu=e^0{}_\mu$ (with $X^0=t$) and the
  spatial cometric $h^{\mu\nu}=\delta^{ij}e_i{}^\mu e_j{}^\nu$.
\item \emph{It defines a connection.} The coefficients
  $\Gamma^\lambda{}_{\mu\nu}=e_a{}^\lambda\partial_\mu e^a{}_\nu$ define the
  derivative for which frame components are constant; they are symmetric in
  $\mu\nu$ (no torsion) and give zero curvature, both consequences of
  $\dd e^a=0$.
\end{enumerate}

\begin{example}[Polar coordinates]\label{ex:polar}
Take the Euclidean plane with $X^1=r\cos\theta$, $X^2=r\sin\theta$. Then
\begin{equation}
  e^1=\cos\theta\,\dd r-r\sin\theta\,\dd\theta,\qquad
  e^2=\sin\theta\,\dd r+r\cos\theta\,\dd\theta,
\end{equation}
so $\delta_{ab}e^a e^b=\dd r^2+r^2\dd\theta^2$, and
$\Gamma^\lambda{}_{\mu\nu}=e_a{}^\lambda\partial_\mu e^a{}_\nu$ gives
$\Gamma^r{}_{\theta\theta}=-r$ and $\Gamma^\theta{}_{r\theta}=1/r$, the familiar
Christoffel symbols of polar coordinates. They are not extra input: they are
what the soldering form of a flat chart looks like in a curvilinear one.
\end{example}

In general relativity the same object appears as the tetrad $e^a{}_\mu$, with
local Lorentz transformations acting on $a$ as a gauge symmetry. The
difference from Kretschmann's construction lies not in the object but in its
status. There $e^a{}_\mu$ is prescribed as the gradient of a given chart, so
$\dd e^a=0$ and the metric $\eta_{ab}e^a{}_\mu e^b{}_\nu$ is flat and the same
in every solution. In general relativity $e^a{}_\mu$ is varied in the action,
$\dd e^a\neq0$ in general (the torsion-free spin connection obeys
$\dd e^a+\omega^a{}_b\wedge e^b=0$), and the metric is dynamical. The
soldering form is the fixed background of every covariantized older theory and
a dynamical field in general relativity, which is the distinction drawn by
Proposition~\ref{prop:main} and Anderson's criterion.

\subsection{Where the physics went}

After the construction the covariance group is $\diff(M)$, acting on pairs
$(\tilde\varphi,b)$ where $b$ is the $G_0$-invariant structure of
Remark~\ref{rem:soldering-dep}. By Proposition~\ref{prop:main} the symmetry
group is $\stab_{\diff(M)}(b)$, and this is exactly the original $G_0$, because
$G_0$ is the automorphism group of its invariant structure. This is Klein's
Erlangen idea read backwards: a geometry determines its group, and the group
is recovered as the automorphism group of an invariant structure, as for
$\mathrm{Poincar\acute e}=\mathrm{Aut}(\eta)$ and
$\text{Galilei}=\mathrm{Aut}(t_\mu,h^{\mu\nu},\nabla)$. Covariantization
therefore changes the \emph{formulation} (the covariance group becomes
$\diff(M)$) and leaves the \emph{physics} (the symmetry group, the solution
set) untouched. Table~\ref{tab:kretschmannize} lists the standard cases. The
friction and preferred-frame rows show what happens when the original theory
contains a genuinely non-symmetric ingredient: the covariantization turns it
into an extra background field, and the stabilizer shrinks accordingly.

\begin{table}[ht]
\centering
\small
\renewcommand{\arraystretch}{1.25}
\setlength{\tabcolsep}{4pt}
\begin{tabular}{@{}>{\raggedright\arraybackslash}p{3.7cm}%
                  >{\raggedright\arraybackslash}p{5.6cm}%
                  >{\raggedright\arraybackslash}p{5.0cm}@{}}
\toprule
Original theory & Absolute structure $b_0$ introduced & Stabilizer $\stab_{\diff(M)}(b_0)$\\
\midrule
Special-relativistic field theory & flat metric $\eta_{\mu\nu}$ & Poincar\'e group\\
Galilean $N$-body mechanics & $t_\mu$, $h^{\mu\nu}$, flat connection $\nabla$ & Galilei group\\
Newtonian gravity (Cartan) & $t_\mu$, $h^{\mu\nu}$; connection dynamical
  & larger than Galilei (includes ${\bm{x}}\to{\bm{x}}+\bm a(t)$)\\
Friction against a medium at rest & as for Galilei, plus the velocity field
  $w^\mu$ of the medium & translations, rotations; no boosts\\
Electrodynamics with a preferred frame & $\eta_{\mu\nu}$ plus a unit timelike
  vector $u^\mu$ & translations, rotations fixing $u^\mu$; no boosts\\
General relativity & none (the metric is dynamical) & $\diff(M)$, all gauge\\
\bottomrule
\end{tabular}
\caption{Kretschmann's construction applied to several theories.}
\label{tab:kretschmannize}
\end{table}

\subsection{Two attempted escapes, and why they do not help}

\paragraph{Make the coordinates dynamical.}
One may try to remove the background by promoting the preferred chart to
dynamical fields. For a Poincar\'e-invariant action $S[\varphi;\eta]$ define the
\emph{parametrized} theory
\begin{equation}
  S_P[\varphi,X]=S\big[\varphi;\ \eta_{\mu\nu}=\eta_{ab}\,\partial_\mu X^a\partial_\nu X^b\big],
\end{equation}
now varied with respect to $X^a$ as well \cite{kuchar}. The action is
$\diff(M)$-invariant, no chart is singled out, and by Noether's second theorem
\eqref{eq:noether2} the equations for $X^a$ are consequences of those for
$\varphi$ (they express the conservation of the energy--momentum tensor). The
$X^a$ are pure gauge: in the gauge $X^a=y^a$ one recovers special relativity.
Nothing physical has changed. By Anderson's criterion this is not a
counterexample, since the flat metric built from $X^a$ is the same in every
solution up to diffeomorphism and is therefore still an absolute object, but
it shows that the criterion must be applied to \emph{composite} objects built
from the dynamical fields as well, and that the mere presence of
diffeomorphism-gauge freedom does not distinguish general relativity from
special relativity \cite{giulini}.

\paragraph{Count the absolute objects.}
Anderson's definition calls an object \emph{absolute} if it appears in the
equations and is the same, up to diffeomorphism, in every solution, and
\emph{dynamical} otherwise. In these terms the distinguishing feature of general
relativity is that the metric is dynamical in the strong sense: different
solutions are not isometric (they differ in their Weyl curvature, and there
are gravitational waves), while every covariantized older theory carries an
inertial structure that is fixed once and for all.

\subsection{What survives of general covariance}

\begin{enumerate}[leftmargin=1.8em,itemsep=3pt,label=\textbf{\arabic*.}]
\item \textbf{As a formal requirement it is vacuous}
  (Proposition~\ref{prop:kretschmann}); it constrains the notation, not the
  theory.
\item \textbf{As gauge invariance it has consequences.} If \emph{all}
  geometric structures are dynamical, $\diff(M)$ is a gauge group and
  Noether's second theorem gives the Bianchi identities. The consequence is
  that $\nabla_\mu T^{\mu\nu}=0$ follows from the gravitational field equations
  alone, for any matter content, so the matter equations must be compatible
  with it. In a theory with a fixed $\eta$ the same conservation law is
  instead a consequence of the matter equations. This is the part of Einstein's
  intuition that the G\"ottingen group made precise \cite{ks}.
\item \textbf{As a heuristic and a uniqueness statement it has content}, as
  Einstein argued in reply. In four dimensions the only symmetric,
  divergence-free tensor built from the metric and its first two derivatives,
  linear in the second derivatives, is $G_{\mu\nu}+\Lambda g_{\mu\nu}$
  \cite{lovelock}: combined with general covariance and a dynamical metric this
  fixes the gravitational field equations almost uniquely.
\item \textbf{The testable content lies in the stabilizer, not in
  covariance.} Universality of free fall, tested to parts in $10^{13}$ and
  beyond, and the null results of Hughes--Drever-type experiments say that
  matter couples to a single metric and that no further background structure
  (a preferred frame $u^\mu$, say) enters the equations. In the language of
  Proposition~\ref{prop:main}: they test \emph{what is held fixed}, 
 covariance, cannot be   tested at all.
\end{enumerate}

\begin{remark}[The Kretschmann scalar]\label{rem:kretschmann-scalar}
The curvature invariant
\begin{equation}
  K=R_{\mu\nu\rho\sigma}R^{\mu\nu\rho\sigma}
\end{equation}
also carries Kretschmann's name. It illustrates the difference between covariant
\emph{components} and invariant \emph{scalars}. For the Schwarzschild solution
\begin{equation}
  K=\frac{48\,G^2M^2}{c^4\,r^6}
\end{equation}
is finite at $r=2GM/c^2$, although $g_{tt}$ vanishes and $g_{rr}$ diverges
there in Schwarzschild coordinates: the horizon is a coordinate artifact,
while $r=0$ is a genuine curvature singularity. Statements about
chart-dependent components are statements about the description; statements
about invariants are statements about the world.
\end{remark}

\section{Example IV: gauge covariance}
\label{sec:gauge}

The vocabulary of Section~\ref{sec:three-words} is stretched, in a
precise way, in gauge theory. Let a charged field $\psi$ with charge $q$
transform under local $U(1)$:
\begin{equation}\label{eq:u1}
  \psi\to e^{-\ii q\Lambda(x)/\hbar c}\,\psi,\qquad
  A_\mu\to A_\mu+\partial_\mu\Lambda .
\end{equation}
The ordinary derivative $\partial_\mu\psi$ is neither invariant nor covariant:
it picks up an extra term $\propto\partial_\mu\Lambda\,\psi$. The
\emph{covariant derivative}
\begin{equation}\label{eq:covder}
  D_\mu=\partial_\mu+\frac{\ii q}{\hbar c}A_\mu
  \quad\Longrightarrow\quad
  D_\mu\psi\to e^{-\ii q\Lambda/\hbar c}\,D_\mu\psi
\end{equation}
transforms with the \emph{same phase} as $\psi$ itself: it is
gauge-\emph{covariant} in the sense of \eqref{eq:covariant-object}, with
$D(g)$ the $U(1)$ phase. Gauge-\emph{invariant} objects are built from
covariant ones by contracting the phases away:
\begin{equation}
  \psi^\dagger\psi,\qquad(D_\mu\psi)^\dagger(D^\mu\psi),\qquad
  F_{\mu\nu}=\frac{\hbar c}{\ii q}[D_\mu,D_\nu]=\partial_\mu A_\nu-\partial_\nu A_\mu .
\end{equation}
For the non-Abelian case, $\psi\to U\psi$, $D_\mu=\partial_\mu-\ii gA_\mu$
with $A_\mu=A_\mu^aT^a$, and
\begin{equation}
  A_\mu\to UA_\mu U^{-1}-\tfrac{\ii}{g}(\partial_\mu U)U^{-1},\quad
  F_{\mu\nu}=\partial_\mu A_\nu-\partial_\nu A_\mu-\ii g[A_\mu,A_\nu]\to UF_{\mu\nu}U^{-1}.
\end{equation}
$F_{\mu\nu}$ is gauge-\emph{covariant}, and $\tr F_{\mu\nu}F^{\mu\nu}$ is
gauge-\emph{invariant}. In the Abelian case $F_{\mu\nu}$ is itself invariant,
not merely covariant, because the adjoint representation of $U(1)$ is trivial.

The interplay with Noether's theorems is instructive. The \emph{global} subgroup
($\Lambda=\text{const}$) is an honest symmetry with a conserved, measurable
charge, the electric charge of Table~\ref{tab:noether}. The \emph{local}
transformations are redundancy: they yield the Noether identity
$\partial_\mu j^\mu\approx0$ with $j^\mu=\partial\mathcal L_{\rm matter}/\partial A_\mu$
(Example~\ref{ex:noether2-em}), and, in the Hamiltonian formalism,
a first-class constraint, Gauss's law $\nabla\!\cdot{\bm{E}}=4\pi\rho$, that
generates the gauge transformations.

\section{Example VI: quantum mechanics}
\label{sec:quantum}

By Wigner's theorem a symmetry of a quantum system, a bijection of rays that
preserves transition probabilities $|\langle\psi|\varphi\rangle|^2$, is
implemented by an operator $U(g)$ that is unitary or antiunitary, unique up to
a phase \cite{wigner}. The phase ambiguity means the group is represented
\emph{projectively}:
$U(g_1)U(g_2)=e^{\ii\omega(g_1,g_2)}U(g_1g_2)$. For rotations this gives
spin-$\tfrac12$ representations of the double cover $SU(2)$; for the
Galilei group the nontrivial phase is the mass $M$ of
Section~\ref{sec:galilei}, the central charge appearing in the classical
Poisson algebra.

The three notions take the following operator form.
\begin{itemize}[leftmargin=1.6em,itemsep=2pt]
\item \emph{Invariance} of the dynamics: $U(g)\,H\,U(g)^{-1}=H$, i.e.\
  $[U(g),H]=0$. Symmetry of the law is invariance of one particular operator.
\item \emph{Covariance} of observables: an observable $\mathcal O$
  transforms into $U^\dagger\mathcal OU$. A \emph{vector operator} $V_i$
  (position, momentum, angular momentum) transforms covariantly,
  \begin{equation}\label{eq:vector-op}
    U(R)^\dagger V_iU(R)=R_{ij}V_j\quad\Longleftrightarrow\quad
    [J_i,V_j]=\ii\hbar\,\epsilon_{ijk}V_k,
  \end{equation}
  whereas the scalars $H$ and $\bm J^2$ are \emph{invariant}.
\item \emph{Consequences.} If $[U,H]=0$, energy levels fall into multiplets
  that carry representations of the symmetry group (the $(2j{+}1)$-fold
  degeneracy for rotations), and matrix elements of covariant operators
  between such multiplets are constrained by the Wigner--Eckart theorem,
  yielding selection rules. Both facts follow from group theory alone,
  without solving the dynamics.
\end{itemize}
The last bullet displays the practical power of the whole subject:
covariance of operators and invariance of $H$ together determine selection
rules before any wavefunction is computed. The classical analogue of
\eqref{eq:vector-op} is $\pb{L_i}{V_j}=\epsilon_{ijk}V_k$.

\begin{remark}[Discrete symmetries]
Parity $P$ and time reversal $T$ are improper transformations. Electrodynamics
is covariant (hence symmetric) under both if $\tilde F$, ${\bm{B}}$ and ${\bm{E}}$ are
assigned the appropriate parity and time-reversal signs; the weak interaction
is not symmetric under $P$, so the ``covariance'' of an equation under $P$
is a statement about the formulation only until one also asks whether the
symmetry is a property of the \emph{law} that nature obeys.
\end{remark}

Table~\ref{tab:summary} places the examples in the framework of
Proposition~\ref{prop:main}.

\begin{table}[ht]
\centering
\footnotesize
\renewcommand{\arraystretch}{1.25}
\setlength{\tabcolsep}{3pt}
\begin{tabular}{@{}>{\raggedright\arraybackslash}p{2.2cm}>{\raggedright\arraybackslash}p{3.1cm}>{\raggedright\arraybackslash}p{2.7cm}>{\raggedright\arraybackslash}p{3.5cm}>{\raggedright\arraybackslash}p{3.1cm}@{}}
\toprule
System & Covariance group (of the formulation)
  & Fixed background & Symmetry $=\stab(b_0)$ & Consequence\\
\midrule
Newtonian $N$-body & Galilei (or $\diff$ in Newton--Cartan form)
  & $t_\mu,\ h^{\mu\nu}$ & Galilei, 10 parameters
  & $E,\bm P,\bm L,\bm G$; mass as central charge\\
Maxwell with sources & Lorentz (or $\diff$ with $\eta$ transported)
  & $\eta_{\mu\nu}$ & Poincar\'e; conformal if $j=0$
  & $P^\mu,J^{\mu\nu}$\\
Relativistic particle & Poincar\'e, reparametrization
  & $\eta_{\mu\nu}$ & Poincar\'e; reparametrization is gauge
  & $\pi_\mu,J^{\mu\nu}$; $\pi^2-m^2c^2\approx0$\\
General relativity & $\diff(M)$ & none & $\diff(M)$, entirely gauge
  & $\nabla_\mu T^{\mu\nu}=0$; charges only for Killing $\xi$ of a solution\\
Gauge theory & local $G$ & $\eta_{\mu\nu}$
  & Poincar\'e $\times$ global $G$; local $G$ is gauge
  & charge; $\partial_\mu j^\mu\approx0$; Gauss law\\
Kepler problem & canonical transformations & --- 
  & $SO(3)\times\mathbb R_t$ (geometric); $SO(4)$ (dynamical, $H<0$)
  & ${\bm{L}},\ \bm A,\ H$\\
Quantum mechanics & unitary / antiunitary operators & fixed by $H$
  & $[U(g),H]=0$
  & degeneracy multiplets; selection rules\\
\bottomrule
\end{tabular}
\caption{The examples of this chapter in the language of
Proposition~\ref{prop:main}.}
\label{tab:summary}
\end{table}

